\section{Conclusion}
We study sycophancy with black-box framework through the lens of pre-pressure uncertainty, asking whether a model's belief state before any pressure begins predicts how it responds to escalating social pressure in a multi-turn dialogue. Across GPT, Claude, and Gemini model families on MMLU-Pro, GPQA-Diamond, and HLE, we find that uncertain questions flip at nearly twice the rate of certain ones. We understand that pressure level amplifies pre-existing uncertainty rather than independently inducing capitulation, and that models divide into resisters and capitulators in ways that do not reduce to model scale, with a third mode, fragmenters, which lose coherent self-consistency under pressure rather than committing to a wrong answer, emerging on HLE. Reasoning traces further expose capitulation earlier than final answer trajectories, with belief entropy spiking at the first pressure turn. These findings reframe sycophancy from a static model property to an interaction between pre-pressure belief state and social pressure structure, and suggest that uncertainty measured through repeated sampling is a model-agnostic signal for identifying fragile questions before any adversarial interaction occurs.
% We have introduced a black-box framework for measuring per-turn semantic entropy under adversarial sycophancy pressure, and demonstrated that early entropy dynamics ($dH/dt$, $T_0 \to T_2$) are a statistically significant and model-agnostic predictor of sycophantic capitulation. The framework reveals distinct per-model entropy signatures and characterises reasoning sycophancy as a categorical cluster shift rather than a gradual degradation, and establishes a $2 \times 3$ dynamic adversary design for cross-family and intra-family evaluation of the corrigibility gap. Our results collectively suggest that multi-turn adversarial pressure exposes vulnerabilities that single-turn benchmarks, and even Constitutional AI training, may not eliminate.
